\documentclass[11pt,oneside,openany]{book}
\usepackage[T1]{fontenc}
\usepackage[utf8]{inputenc}
\usepackage[sc]{mathpazo}
\renewcommand{\sfdefault}{lmss}
\renewcommand{\ttdefault}{lmtt}
\usepackage{amsmath,amssymb,amsthm,mathtools}
\usepackage[paperwidth=7.5in,paperheight=10in,inner=.85in,outer=.85in,top=.85in,bottom=.85in,headheight=15pt,headsep=20pt]{geometry}
\usepackage{microtype}
\usepackage{xcolor,graphicx,booktabs,array,longtable}
\newcolumntype{L}[1]{>{\raggedright\arraybackslash}p{#1}}
\usepackage{tikz,pgfplots}
\usetikzlibrary{arrows.meta,calc,patterns,decorations.markings}
\pgfplotsset{compat=1.18}
\definecolor{ink}{HTML}{17334B}
\definecolor{teal}{HTML}{137C83}
\definecolor{copper}{HTML}{AE5932}
\definecolor{soft}{HTML}{EEF3F5}
\pgfplotsset{every axis/.append style={axis line style={black!65},tick style={black!60},tick label style={font=\small},label style={font=\small},legend style={font=\small,draw=none,fill=white},grid style={black!10},major grid style={black!10},line width=.6pt},every axis plot/.append style={line width=1.15pt}}
\usepackage{enumitem,needspace,etoolbox}
\setlist[enumerate]{topsep=5pt,itemsep=5pt,parsep=0pt}
\setlist[itemize]{topsep=4pt,itemsep=3pt,parsep=0pt}
\usepackage{titlesec,fancyhdr}
\titleformat{\chapter}[display]{\normalfont\color{ink}}{\large\sffamily\MakeUppercase{Chapter \thechapter}}{12pt}{\huge\bfseries\raggedright\hyphenpenalty=10000\exhyphenpenalty=10000}[\vspace{12pt}\titlerule]
\titlespacing*{\chapter}{0pt}{-10pt}{24pt}
\titleformat{\section}{\large\bfseries\color{ink}}{\thesection}{.65em}{}
\titleformat{\subsection}{\normalsize\bfseries\color{ink}}{\thesubsection}{.6em}{}
\titlespacing*{\section}{0pt}{18pt plus 3pt minus 2pt}{7pt}
\titlespacing*{\subsection}{0pt}{12pt plus 2pt minus 2pt}{5pt}
\pagestyle{fancy}
\fancyhf{}
\fancyhead[L]{\small\nouppercase{\leftmark}}
\fancyfoot[C]{\small\thepage}
\renewcommand{\headrulewidth}{.3pt}
\renewcommand{\chaptermark}[1]{\markboth{\thechapter\quad #1}{}}
\fancypagestyle{plain}{\fancyhf{}\fancyfoot[C]{\small\thepage}\renewcommand{\headrulewidth}{0pt}}
\usepackage[font=small,labelfont=bf,labelsep=period]{caption}
\usepackage{float,placeins}
\floatplacement{figure}{htbp}
\setcounter{topnumber}{3}
\setcounter{bottomnumber}{2}
\renewcommand{\topfraction}{.9}
\renewcommand{\textfraction}{.08}
\renewcommand{\floatpagefraction}{.8}
\usepackage{makeidx}
\makeindex
\usepackage[pdfencoding=auto,psdextra,colorlinks=true,linkcolor=ink,citecolor=teal,urlcolor=teal,bookmarksnumbered=true]{hyperref}
\usepackage{bookmark}
\hypersetup{pdftitle={Discrete Information Geometry},pdfsubject={Finite differences, probability, and information flows}}
\urlstyle{same}
\newtheorem{theorem}{Theorem}[chapter]
\newtheorem{proposition}[theorem]{Proposition}
\newtheorem{lemma}[theorem]{Lemma}
\newtheorem{corollary}[theorem]{Corollary}
\theoremstyle{definition}
\newtheorem{definition}[theorem]{Definition}
\newtheorem{example}{Example}[chapter]
\BeforeBeginEnvironment{example}{\par\addvspace{4pt}\noindent\begin{minipage}{\linewidth}}
\AfterEndEnvironment{example}{\end{minipage}\par\addvspace{4pt}}
\theoremstyle{remark}
\newtheorem*{remark}{Remark}
\newenvironment{solution}{\par\smallskip\noindent\textit{Solution.}\ }{\par\medskip}
\newenvironment{exercises}{\begin{enumerate}[label=\thechapter.\arabic*,ref=\thechapter.\arabic*,leftmargin=2.8em,itemsep=6pt]}{\end{enumerate}}
\newenvironment{answers}{\begin{description}[leftmargin=2.9em,labelwidth=2.5em,labelsep=.4em,itemsep=8pt,font=\normalfont\bfseries]}{\end{description}}
\newcommand{\R}{\mathbb{R}}
\newcommand{\E}{\mathbb{E}}
\newcommand{\Var}{\operatorname{Var}}
\newcommand{\Cov}{\operatorname{Cov}}
\newcommand{\KL}[2]{D\!\left(#1\,\middle\|\,#2\right)}
\newcommand{\grad}{\operatorname{grad}}
\newcommand{\arcosh}{\operatorname{arcosh}}
\newcommand{\sech}{\operatorname{sech}}
\newcommand{\softmax}{\operatorname{softmax}}
\newcommand{\diag}{\operatorname{diag}}
\newcommand{\dd}{\,\mathrm{d}}
\newcommand{\ip}[2]{\left\langle #1,#2\right\rangle}
\newcommand{\p}{\mathbf{p}}
\newcommand{\q}{\mathbf{q}}
\newcommand{\f}{\mathbf{f}}
\newcommand{\uvec}{\mathbf{u}}
\newcommand{\vvec}{\mathbf{v}}
\newcommand{\ones}{\mathbf{1}}
\newcommand{\simp}{\Delta^{n-1}}
\newcommand{\simpi}{\Delta^{n-1}_{\!\circ}}
\numberwithin{equation}{chapter}
\setlength{\parindent}{1.2em}
\setlength{\parskip}{0pt}
\linespread{1.07}
\setlength{\emergencystretch}{2em}
\clubpenalty=10000
\widowpenalty=10000
\displaywidowpenalty=10000
\raggedbottom
\setcounter{tocdepth}{1}
\makeatletter
\renewcommand*\l@section{\@dottedtocline{1}{1.5em}{3.2em}}
\makeatother
\setcounter{secnumdepth}{2}
\providecommand{\tightlist}{\setlength{\itemsep}{0pt}\setlength{\parskip}{0pt}}
\providecommand{\passthrough}[1]{#1}
\newcommand{\labheading}{\FloatBarrier\section*{Laboratory investigation}\addcontentsline{toc}{section}{Laboratory investigation}}

\begin{document}
\begin{titlepage}\color{ink}\centering\vspace*{1.0in}{\small\sffamily PROBABILITY AND MATHEMATICAL MODELLING\par}\vspace{.45in}{\Huge\bfseries Discrete Information\\[.12in]Geometry\par}\vspace{.35in}{\Large Finite Differences, Probability,\\and Information Flows\par}\vspace{1cm}{\Large Chapter 9\par}\vfill\begin{tikzpicture}[scale=1.05]\draw[teal,line width=1pt](0,0)--(4,0)--(2,3.464)--cycle;\foreach \a in {0,...,6}{\foreach \b in {0,...,6}{\pgfmathtruncatemacro{\c}{6-\a-\b}\ifnum\c>-1\fill[teal] ({(\b+\c/2)*4/6},{\c*3.464/6}) circle(1.5pt);\fi}}\end{tikzpicture}\vfill{\large A course in ten chapters\par}\vspace{.3in}{\small\url{https://bobovski66.github.io/IntroDiscreteInfoGeom/}\par}\end{titlepage}
\setcounter{chapter}{8}

\chapter{Discrete natural-gradient ideas and learning dynamics}
Suppose an algorithm divides attention among three prediction methods. Each method receives a score after a round of observations. The algorithm should give more attention to better-performing methods, while retaining a valid probability distribution and avoiding an unnecessarily abrupt change.

Earlier chapters introduced all the ingredients for such a rule. Scores specify a preferred direction. Relative entropy measures the cost of departure from the current table. Fisher geometry describes the scale of small changes. We will combine these ideas into an exact finite optimization problem, solve it algebraically, and then explain its connection to natural gradients.

\section{A score specifies what an update should improve}\label{ch09-a-score-specifies-what-an-update-should-improve}

Let \(p=(p_1,\ldots,p_k)\) be current probabilities, initially with every \(p_i>0\). Let \(s_i\) be the score of choice \(i\). If a proposed new table is \(q\), its average score is \[q\cdot s=\sum_i q_i s_i.\] This is a linear function of \(q\). Without any restriction on change, maximizing it would place all probability on a highest-scoring choice.

That may be too abrupt. Scores may be noisy, and the current table may encode accumulated evidence. Introduce a positive intensity \(\eta\) and maximize \[\boxed{J(q)=\eta\,q\cdot s-D(q\|p)}\] over probability vectors \(q\).

The score term rewards improvement; the relative-entropy term charges for changing the current table. Larger \(\eta\) gives scores greater importance relative to that charge. The product \(\eta s_i\) must be dimensionless, so score units and intensity should be chosen consistently.

The direction of KL is part of this objective: the candidate \(q\) supplies the averaging weights, and the old state \(p\) is the reference. Reversing those arguments defines a different problem.

A score can represent a reward or the negative of a loss. If \(s_i\) is an error count, directly favoring larger \(s_i\) would reward worse predictions. To discourage loss \(\ell_i\), use \(s_i=-\ell_i\), giving multipliers \(e^{-\eta\ell_i}\). All the algebra remains valid; the interpretation is determined by which values are desirable.

The current probabilities matter even when scores are fixed. They can represent past evidence, an initial allocation, or the frequency with which choices are used. The update is therefore a response to both the prior state and the new score, rather than a fresh ranking that discards history.

A tie in scores illustrates this clearly. If \(s_i=s_j\), then \(r_i/r_j=p_i/p_j\): their relative weight is unchanged. If every score is equal, the whole table remains unchanged. This provides a quick consistency check on both the formula and a numerical implementation.

The objective also explains the role of the intensity. At very small intensity, the departure penalty dominates and the optimizer remains close to \(p\). At a large intensity, distinctions among scores matter more. Exact probabilities remain nonnegative at every finite intensity, but numerical underflow can make a very small weight appear zero. That computational artifact should be distinguished from the support behavior of the exact recurrence.

\section{Reweight, then normalize}\label{ch09-reweight-then-normalize}

The optimizer will be \[\boxed{r_i=\frac{p_i e^{\eta s_i}}{Z},\qquad
Z=\sum_j p_j e^{\eta s_j}.}\] First multiply each current probability by a positive exponential score factor. Then normalize the resulting weights. This is called an \textbf{exponential update} or a \textbf{multiplicative-weights update}.

Every new entry is positive and the total is one. The rule is the same normalized multiplication used in population selection, now interpreted as a learning step.

\begin{example}[Redistributing attention among three methods]

Take \(p=(1/2,1/3,1/6)\), scores \(s=(0,1,2)\), and \(\eta=\ln2\). The multipliers are \((1,2,4)\). Thus \[Z=\frac12+\frac23+\frac23=\frac{11}6,\qquad
r=\left(\frac3{11},\frac4{11},\frac4{11}\right).\] The third method has the highest score but starts with the smallest probability. Its updated probability ties the second method because its larger multiplier compensates for its smaller starting weight.

\end{example}

\begin{figure}[htbp]\centering
\resizebox{.98\linewidth}{!}{\begin{tikzpicture}\begin{axis}[width=12.5cm,height=5.6cm,ybar,bar width=14pt,xlabel={Choice},ylabel={Probability},xmin=.5,xmax=3.5,ymin=0,ymax=.65,xtick={1,2,3},legend pos=north east]
\addplot[ink,fill=ink!40] coordinates {(1,0.5) (2,0.333333) (3,0.166667)};\addlegendentry{Current table}
\addplot[teal,fill=teal!60] coordinates {(1,0.272727) (2,0.363636) (3,0.363636)};\addlegendentry{Exponential update}
\end{axis}\end{tikzpicture}}
\caption{An exponential update of (1/2,1/3,1/6) with scores (0,1,2) and intensity ln 2 gives (3/11,4/11,4/11).}\end{figure}

Subtracting the old table gives the discrete derivative \[\Delta p=r-p.\] For the example, its coordinates are \((-5/22,1/33,13/66)\), whose sum is zero. The negative first coordinate records attention released to the other choices.

\section{Why this table solves the optimization problem}\label{ch09-why-this-table-solves-the-optimization-problem}

The proof follows by expanding one relative entropy. From the formula for \(r_i\), \[\ln r_i=\ln p_i+\eta s_i-\ln Z.\] Substitute into \(D(q\|r)\): \[\begin{aligned}
D(q\|r)
&=\sum_i q_i(\ln q_i-\ln r_i)\\
&=D(q\|p)-\eta\sum_iq_i s_i+\ln Z.
\end{aligned}\] The last term is \(\ln Z\) because \(\sum_iq_i=1\). Rearranging gives \[\boxed{J(q)=\ln Z-D(q\|r).}\] Gibbs' inequality makes the right side at most \(\ln Z\), with equality exactly when \(q=r\). Therefore \(r\) is the unique optimizer.

The identity does more than identify a best table. It measures the loss from any candidate: \[J(r)-J(q)=D(q\|r).\] A candidate's directed discrepancy from the optimizer is exactly its objective shortfall.

\begin{example}[Comparing an unchanged state with the optimizer]

For \(p=(1/2,1/2)\), scores \((0,1)\), and \(\eta=\ln3\), the update is \(r=(1/4,3/4)\) and \(Z=2\). The maximum objective is \(\ln2\). At the unchanged table, \[J(p)=\frac12\ln3.\] Their difference is \(\ln2-\tfrac12\ln3=\tfrac12\ln(4/3)\), equal to \(D(p\|r)\).

\end{example}

The optimizer ranges over all probability vectors, including boundary candidates. The problem concerns a finite-step update on a finite label set; the set of possible real-valued probability vectors is still continuous. If candidates must belong to a token lattice, the exact optimizer may be unavailable, and the identity can instead compare the permitted candidates.

\section{What the update guarantees about scores}\label{ch09-what-the-update-guarantees-about-scores}

Because \(p\) itself is an available candidate and \(D(p\|p)=0\), \[\eta\,r\cdot s-D(r\|p)\geq\eta\,p\cdot s.\] Hence \[\boxed{\eta(r-p)\cdot s\geq D(r\|p)\geq0.}\] The new table has at least as large an average under the scores used to choose the update. When those scores differ on the positive support, the table changes, the divergence is positive, and the average strictly improves.

This is a statement about a \textbf{fixed score vector during the step}. If scores are recomputed from the new state, the next vector of scores may differ. Improving the old average does not automatically improve an average evaluated under new scores. Cyclic interactions from Chapter 3 are a useful reminder of this distinction.

For the three-method example, the average rises from \(2/3\) to \(12/11\). The one-step improvement is \(14/33\). Repeatedly using the same scores keeps favoring the same methods. A learning problem with new observations at each round generally supplies changing scores.

\section{Odds reveal an exact additive update}\label{ch09-odds-reveal-an-exact-additive-update}

For two positive coordinates, the common normalization cancels: \[\frac{r_i}{r_j}=\frac{p_i}{p_j}e^{\eta(s_i-s_j)}.\] Taking logarithms yields \[\boxed{\ln\frac{r_i}{r_j}-\ln\frac{p_i}{p_j}
=\eta(s_i-s_j).}\] Relative log weights have an exact finite difference equal to the score gap times the intensity.

For binary choices with \(p=(x,1-x)\) and \(d=s_A-s_B\), \[x^+=\frac{xe^{\eta d}}{xe^{\eta d}+1-x}.\] At \(x=1/2\), \(d=\ln2\), and \(\eta=1\), the next value is \(2/3\). The proportion changes by \(1/6\), while the log odds change by \(\ln2\).

This is the odds simplification from Chapter 1 expressed on a logarithmic scale. Repeated multiplication has become accumulated addition.

\section{The geometry determines a local direction}\label{ch09-the-geometry-determines-a-local-direction}

Why is the update connected with a \textbf{natural gradient}? An ordinary gradient chooses steepness using ordinary Euclidean length. A natural gradient uses the statistical length supplied by the Fisher metric. We can derive the direction by completing a square, without differentiating a function.

Consider a small displacement \(\delta\) with \(\sum_i\delta_i=0\). Approximate the change penalty using the Fisher form from Chapter 7: \[D(p+\delta\|p)\approx\frac12\sum_i\frac{\delta_i^2}{p_i}.\] The part of the objective that depends on the displacement is approximately \[Q(\delta)=\eta\sum_i\delta_i s_i-\frac12\sum_i\frac{\delta_i^2}{p_i}.\]

Let \(\bar s=\sum_i p_i s_i\). Since \(\sum_i\delta_i=0\), replacing \(s_i\) by \(s_i-\bar s\) does not change the linear term. Completing the square gives \[\begin{aligned}
Q(\delta)
&=\frac{\eta^2}{2}\sum_i p_i(s_i-\bar s)^2\\
&\quad-\frac12\sum_i
\frac{[\delta_i-\eta p_i(s_i-\bar s)]^2}{p_i}.
\end{aligned}\] The second sum is nonnegative. It is zero exactly at \[\boxed{\delta_i=\eta p_i(s_i-\bar s).}\] These coordinates add to zero, so the proposed displacement satisfies the mass constraint.

Thus \[v_i=p_i(s_i-\bar s)\] is the Fisher natural-gradient direction for a linear score function. The factor \(p_i\) reflects the local statistical geometry. The subtraction of the mean score enforces conservation of total probability.

A large step \(p+\eta v\) may acquire negative entries. This local quadratic calculation therefore does not replace the exact finite update. It explains its small-step direction.

The same direction can be characterized by a fixed statistical movement budget. Suppose \(\sum_i\delta_i=0\) and \(g_p(\delta,\delta)\leq\varepsilon^2\). The score improvement is \[\sum_i\delta_i s_i
=\sum_i\frac{\delta_i}{\sqrt{p_i}}\sqrt{p_i}(s_i-\bar s).\] Cauchy--Schwarz bounds it by \[\varepsilon\sqrt{\sum_i p_i(s_i-\bar s)^2}.\] Equality is attained when \(\delta_i\) is a positive multiple of \(p_i(s_i-\bar s)\), unless all scores are equal. The sum-zero constraint is satisfied by that vector.

Thus the Fisher natural gradient identifies the direction of greatest score increase for a specified small statistical budget. The completed-square calculation chooses a step along that direction by trading gain against a quadratic cost. The exact exponential optimization uses the full KL cost instead of its local quadratic approximation.

This interpretation makes the geometry operational. A rare coordinate is expensive to change by a fixed absolute amount because its squared displacement is divided by a small \(p_i\). The maximizing direction adjusts changes to that cost while also preserving total probability. Steepness is meaningful only after the measure of movement has been chosen.

\section{Comparing the finite update with the local direction}\label{ch09-comparing-the-finite-update-with-the-local-direction}

For small \(\eta\), the exponential satisfies \(e^{\eta s_i}\approx1+\eta s_i\). Its normalizing factor is approximately \(1+\eta\bar s\). Dividing and retaining first-order changes gives \[r_i-p_i\approx\eta p_i(s_i-\bar s).\] The same direction emerges from the exact update and from the local quadratic optimization.

For \(p=(1/2,1/3,1/6)\) and \(s=(0,1,2)\), the mean is \(2/3\) and \[v=\left(-\frac13,\frac19,\frac29\right).\] For a small intensity, the change should be close to \(\eta v\). As intensity grows, the exact normalization accounts for finite effects that this first-order approximation omits.

The relation is also the algebra behind the \textbf{replicator direction} in population dynamics: a type gains share when its score exceeds the population average. The present course studies the exact finite recurrence. The companion calculus course develops the limiting differential equation.

\section{Intensity, support, and stable computation}\label{ch09-intensity-support-and-stable-computation}

Adding a common number \(C\) to all scores multiplies every unnormalized weight by \(e^{\eta C}\), which cancels. Multiplying all scores by a positive number rescales the intensity. Only score differences affect the normalized result.

Let \(s_{\min}\) and \(s_{\max}\) be the smallest and largest scores, and let \(R=s_{\max}-s_{\min}\). Since \(Z\) is a weighted average of the exponentials, \[e^{\eta s_{\min}}\leq Z\leq e^{\eta s_{\max}}.\] It follows that \[e^{-\eta R}\leq\frac{r_i}{p_i}\leq e^{\eta R}.\] This gives a finite bound on relative reweighting. A small product \(\eta R\) limits how dramatically one step can change each positive coordinate.

For computation, subtract \(s_{\max}\) before exponentiating. The update becomes \[r_i=\frac{p_i e^{\eta(s_i-s_{\max})}}
{\sum_jp_j e^{\eta(s_j-s_{\max})}}.\] The largest exponent is now zero, reducing overflow risk without changing the answer.

If a current coordinate is zero, multiplicative updating keeps it zero. The optimization explanation agrees: a candidate assigning positive mass there would have infinite \(D(q\|p)\). To reopen missing choices, one can first mix with a positive reference table. That is an additional operation with a clear interpretation.

A finite update should be evaluated at the scale at which its assumptions apply. If a score summarizes one new batch of observations, using that score repeatedly without receiving new data would count the same evidence more than once in a Bayesian interpretation. In a fixed-payoff population model, repeated use may be exactly what the model intends. The formula alone does not settle this distinction.

It is also useful to separate the score used for an update from the score used to evaluate its eventual performance. The one-step improvement inequality is exact for the supplied vector. It does not guarantee that a prediction method will perform better on the next unseen batch. Such a claim would need assumptions relating future observations to current scores.

When comparing update intensities experimentally, keep the score vector and current distribution fixed at first. Compare \(\Delta p\), the new average score, and \(D(r\|p)\). These measurements expose how a larger response trades greater score improvement against greater movement. After that controlled comparison, allow scores to vary across rounds.

\section{Repeated rounds accumulate evidence}\label{ch09-repeated-rounds-accumulate-evidence}

For a specified score sequence \(s_{i,n}\) and constant intensity, \[p_{i,T}=
\frac{p_{i,0}\exp\left(\eta\sum_{n=0}^{T-1}s_{i,n}\right)}
{\sum_jp_{j,0}\exp\left(\eta\sum_{n=0}^{T-1}s_{j,n}\right)}.\] Each round multiplies the previous weights by another exponential. All intermediate normalizing factors cancel in the final ratio.

If scores are fixed at \(c_i\), this reduces to \[p_{i,T}=\frac{p_{i,0}e^{T\eta c_i}}
{\sum_jp_{j,0}e^{T\eta c_j}}.\] A uniquely highest-scoring available choice approaches probability one. Tied leaders retain their initial ratios while lower-scoring choices fade. Finite positive multipliers preserve positive coordinates at every finite time, even when their limiting values are zero.

If score vectors are supplied independently of the evolving state, their order does not change the final accumulated scores. Intermediate states can still differ. When scores depend on the state, reordering rounds may change the scores themselves, so this order-independence conclusion no longer follows.

\section{Bayesian updating is a particular reweighting rule}\label{ch09-bayesian-updating-is-a-particular-reweighting-rule}

Suppose \(p_i\) is a current probability for hypothesis \(i\), and an observation has likelihood \(L_i>0\) under that hypothesis. Bayes' rule gives \[r_i=\frac{p_iL_i}{\sum_jp_jL_j}.\] Setting \(s_i=\ln L_i\) and \(\eta=1\) turns this into the exponential update. Likelihood is the probability of the observation under a hypothesis; it is not the hypothesis's prior probability.

Start with equal probabilities for three hypotheses. One observation has likelihoods \((0.8,0.4,0.4)\), proportional to \((2,1,1)\). A second conditionally independent observation has likelihoods \((0.25,0.75,0.25)\), proportional to \((1,3,1)\). Multiplying gives weights proportional to \((2,3,1)\), hence posterior probabilities \((1/3,1/2,1/6)\).

The example connects relative growth, accumulated log evidence, and learning within the same normalization formula. The interpretation depends on what the weights and scores represent.

Conditional independence is the assumption that allows the likelihoods of the two observations to multiply once a hypothesis is fixed. Without it, use the second observation's conditional likelihood given both the hypothesis and the first observation. Bayes' rule still reweights, but the table of multipliers may change.

A zero likelihood has a direct meaning: the hypothesis says the observed event is impossible. It removes that hypothesis from the posterior support. Our exponential formula with finite scores uses positive likelihoods because \(\ln0\) is not finite. A model that permits exact zeros should handle them explicitly, rather than substitute an arbitrary very negative score and call the result exact.

The optimization interpretation and the evidence interpretation therefore overlap under stated assumptions. Both normalize positive reweighting factors, while each supplies its own reason for choosing those factors. This is why the recurrence appears in several subjects without making their underlying mechanisms identical.

\section{Seeking uncertainty is another objective}\label{ch09-seeking-uncertainty-is-another-objective}

A high-score update need not increase entropy. If the scores favor one option, the table may become more concentrated. To seek uncertainty deliberately, a simple rule is mixing toward uniform: \[p^+=(1-\lambda)p+\lambda u.\] Concavity and maximal uniform entropy give \[H(p^+)\geq(1-\lambda)H(p)+\lambda H(u)\geq H(p).\] Relative entropy to uniform decreases at the same time because it is the entropy gap.

The \href{https://bobovski66.github.io/IntroDiscreteInfoGeom/demos/week09.html}{learning demonstration} lets you compare finite reweighting steps. Begin with the worked three-method example, record the displacement, and then lower the intensity to compare \((r-p)/\eta\) with the natural-gradient direction.

We can now choose an update by an objective and explain its geometry. Chapter 10 asks what remains observable when distributions move through a network or pass through a noisy recording system.

\section{Exercises}\label{ch09-exercises}

\begin{exercises}

\item 

Compute the exponential update of \((1/4,1/4,1/2)\) with scores \((2,1,0)\) and \(\eta=\ln2\).

\item 

Derive \(J(q)=\ln Z-D(q\|r)\) and identify where normalization of \(q\) enters the proof.

\item 

For \(p=(1/2,1/2)\), scores \((0,1)\), and \(\eta=\ln3\), compare the objectives at \(p\), the optimizer, and \((0,1)\).

\item 

Use optimality to prove \(\eta(r-p)\cdot s\geq D(r\|p)\).

\item 

Derive the exact log-ratio increment and use it for two rounds with binary score gaps \(\ln2\) and \(\ln3\), starting at equal probabilities and \(\eta=1\).

\item 

Calculate the natural-gradient direction for \(p=(1/2,1/3,1/6)\) and scores \((0,1,2)\). Verify its coordinates sum to zero.

\item 

Expand the completed-square formula for \(Q(\delta)\) and explain why its maximizer satisfies the mass constraint.

\item 

Starting from equal probabilities, use fixed scores \((0,1,1)\) and \(\eta=\ln2\). Find the table after four rounds.

\item 

Repeat Exercise 8 from \((1/2,1/3,1/6)\) and explain the ratio of the tied leaders.

\item 

Show that shifting scores \((1000,1001,999)\) by their largest entry leaves the normalized answer unchanged.

\item 

Apply the two Bayesian likelihood updates in the worked example. State the independence assumption needed to multiply their likelihoods.

\item 

Compare entropy change under a score-favoring update and under mixing toward uniform. Explain why the update objective matters.

\end{exercises}

\section{Selected solutions}\label{ch09-selected-solutions}

\textbf{1.} The unnormalized weights are \((1,1/2,1/2)\), giving \((1/2,1/4,1/4)\).

\textbf{3.} The three objective values are \(\tfrac12\ln3\), \(\ln2\), and \(\ln(3/2)\), respectively. The optimizer has the largest value.

\textbf{5.} The odds are multiplied first by two and then by three. Final odds are six, giving probabilities \((6/7,1/7)\). The accumulated log-odds change is \(\ln6\).

\textbf{6.} The mean score is \(2/3\), giving \((-1/3,1/9,2/9)\). Their sum is zero.

\textbf{8.} The weights are proportional to \((1,16,16)\), so the table is \((1/33,16/33,16/33)\).

\textbf{9.} Unnormalized weights are \((1/2,16/3,8/3)\), totaling \(17/2\). The table is \((1/17,32/51,16/51)\); the tied leaders retain their original ratio two to one.
\labheading\begin{enumerate}
\def\labelenumi{\arabic{enumi}.}
\tightlist
\item
  Let \(p=(1/2,1/3,1/6)\), scores \(s=(0,1,2)\), and intensity \(\eta=\ln2\). (a) Find the three unnormalized weights. (b) Normalize. (c) Compute each coordinate of \(\Delta p\). (d) Verify that the three differences sum to zero.
\item
  Take \(p=(1/2,1/2)\), score difference \(d=\ln2\), and \(\eta=1\). (a) Calculate the next A probability. (b) Calculate \(\Delta x\). (c) Find the initial and new odds. (d) Verify \(\Delta\log O=\eta d\).
\item
  Using your three-category example, compute \(J(q)=\eta q\cdot s-D(q\|p)\) for the original p and for the exponential optimizer r. (a) Evaluate both. (b) Find \(Z\). (c) Check \(J(r)=\ln Z\).
\item
  Repeat the three-category update for \(\eta=0.1\) and \(\eta=0.5\). (a) Compute the actual coordinate changes. (b) Find \(\eta p_i(s_i-\bar s)\) at each intensity. (c) Compare approximation errors and explain why the smaller step agrees better.
\item
  Design scores for which one category begins rare but becomes more likely. (a) Choose a valid positive starting table. (b) Calculate four updates. (c) List the probability differences and the score-weighted means. (d) Explain what each number reveals about the dynamics.
\end{enumerate}

\FloatBarrier

\end{document}
